% part: intuitionistic-logic
% chapter: propositions-as-types
% section: sequent-natural-deduction

\documentclass[../../../include/open-logic-section]{subfiles}

\begin{document}

\olfileid{int}{pty}{snd}

\olsection{سېکوېنټي فطري استنتاج}

که داسې فطري استنتاجي !!{derivation} وي چې
!!{undischarged} فرضونه ئې $\Gamma$ او پايله ئې
$!A$ وي، نو $\Gamma \Sequent !A$ ليکو؛ که
$\Gamma$ تش وي، $\Sequent !A$ ليکو۔

د $\Gamma \cup \{!A_1, \dots, !A_n\}$ دپاره
$\Gamma, !A_1, \dots, !A_n$، او د
$\Gamma \cup \Delta$ دپاره $\Gamma, \Delta$ ليکو۔

پام وکړئ: که $\Gamma \Sequent !A \land !A$
ولرو، معنا ئې دا ده چې !!a{derivation} شته چې
!!{undischarged} فرضونه ئې $\Gamma$ او وروستے
!!{formula} ئې $!A \land !A$ دے۔ په پاې کښې د
$\Elim{\land}$ په کارولو سره !!a{derivation}
اخيستلے شو چې !!{undischarged} فرضونه ئې عين
وي او پايله ئې $!A$ وي۔ په بله وينا، که
$\Gamma \Sequent !A \land !B$ وي، نو
$\Gamma \Sequent !A$ هم وي۔
\begin{prooftree}
  \Axiom$\Gamma \fCenter !A \land !B$
  \RightLabel{$\Elim{\land}$}
  \UnaryInf$\Gamma \fCenter !A$
  \DisplayProof\qquad\bottomAlignProof
  \Axiom$\Gamma \fCenter !A \land !B$
  \RightLabel{$\Elim{\land}$}
  \UnaryInf$\Gamma \fCenter !B$
\end{prooftree}
د $\Elim{\land}$ نښه په فطري استنتاج کښې د
همدغه نوم له قاعدې سره تړاو ښيي۔

په همدې ډول، فرض کړئ $\Gamma, !A \Sequent !B$
لرو؛ يعنې !!a{derivation} شته چې
!!{undischarged} فرضونه ئې $\Gamma, !A$ او وروستے
!!{formula} ئې~$!B$ دے۔ که د $\Intro{\lif}$ قاعده
وکاروو، !!a{derivation} ترلاسه کېږي چې
!!{undischarged} فرضونه ئې $\Gamma$ او وروستے
!!{formula} ئې $!A \lif !B$ وي؛ يعنې
$\Gamma \Sequent !A \lif !B$۔ پام وکړئ چې د
فرضونو !!{discharge} ئې څومره څرګند کړل۔
\begin{prooftree}
  \Axiom $\Gamma, !A \fCenter !B$
  \RightLabel{$\Intro{\lif}$}
  \UnaryInf $\Gamma \fCenter !A \lif !B$
\end{prooftree}

په همدې طريقه له نورو قواعدو هم پايلې راځي؛
تفصيل ئې دا دے:
\begin{gather*}
  \Axiom $\Gamma \fCenter !A$
  \Axiom $\Delta \fCenter !B$
  \RightLabel{$\Intro{\land}$}
  \BinaryInf $\Gamma,\Delta \fCenter !A \land !B$
  \DisplayProof\\
  \Axiom $\Gamma \fCenter !A \land !B$
  \RightLabel{$\Elim{\land}_1$}
  \UnaryInf $\Gamma \fCenter !A$
  \DisplayProof
  \qquad
  \Axiom $\Gamma \fCenter !A \land !B$
  \RightLabel{$\Elim{\land}_2$}
  \UnaryInf $\Gamma \fCenter !B$
  \DisplayProof
  \\
  \Axiom $\Gamma \fCenter !A$
  \RightLabel{$\Intro{\lor}_1$}
  \UnaryInf $\Gamma \fCenter !A \lor !B$
  \DisplayProof\qquad
  \Axiom $\Gamma \fCenter !B$
  \RightLabel{$\Intro{\lor}_2$}
  \UnaryInf $\Gamma \fCenter !A \lor !B$
  \DisplayProof\\
  \Axiom $\Gamma \fCenter !A \lor !B$
  \Axiom $\Delta, !A \fCenter !C$
  \Axiom $\Delta', !B \fCenter !C$
  \RightLabel{$\Elim{\lor}$}
  \TrinaryInf $\Gamma, \Delta, \Delta' \fCenter !C$
  \DisplayProof
  \\
  \Axiom $\Gamma, !A \fCenter !B$
  \RightLabel{$\Intro{\lif}$}
  \UnaryInf $\Gamma \fCenter !A \lif !B$
  \DisplayProof
  \qquad
  \Axiom $\Delta \fCenter !A \lif !B$
  \Axiom $\Gamma \fCenter !A$
  \RightLabel{$\Elim{\lif}$}
  \BinaryInf $\Gamma, \Delta \fCenter !B$
  \DisplayProof
  \\
  \Axiom $\Gamma \fCenter \lfalse$
  \RightLabel{$\FalseInt$}
  \UnaryInf $\Gamma \fCenter !A$
  \DisplayProof
\end{gather*}

هر فرض په خپله له~$!A$ څخه د $!A$
!!a{derivation} دے؛ يعنې تل $!A \Sequent !A$ لرو۔

\begin{prooftree}
  \AxiomC{$ $}
  \UnaryInfC{$!A \fCenter !A$}
\end{prooftree}

دغه قواعد يوځاے د هغو فطري استنتاجي
!!{derivation}s د شتون حساب ګڼلے شو چې شته۔ د
فطري استنتاج د نښو يوه بله بڼه ئې هم ګڼلے شو،
چې هر ګام ئې يوازې !!{derive}d شوے !!{formula}
نۀ، بلکې هغه !!{undischarged} فرضونه هم ثبتوي
چې ترې !!{derive}d شوے دے۔

\begin{prooftree}
  \Axiom$!A \fCenter !A$
  \UnaryInf $!A \fCenter !A \lor (!A \lif \lfalse)$
  \Axiom $!B \fCenter !B$
  \BinaryInf$!A, !B \fCenter \lfalse$
  \UnaryInf$!B \fCenter !A \lif \lfalse$
  \UnaryInf$!B \fCenter !A \lor (!A \lif \lfalse)$
  \Axiom$!B \fCenter !B$
  \BinaryInf$!B \fCenter \lfalse$
  \UnaryInf$\fCenter !B \lif \lfalse$
\end{prooftree}
دلته $!B$ د $(!A \lor (!A \lif \lfalse))\lif \lfalse$
لنډ نوم دے۔ \textbf{د ژباړې د نسخې سمون (OLCMP-218):}
د سرچينې په وروستي رسم کښې د سياق تر دويم فرض
پسې يوه بې‌دويم آرګومېنټه د شرط نښه، او د
همدغه فرض مخې ته څلور بې‌جوړې پرانيستې قوسونه
وو۔ يوازې دغه پنځه نقل شوې نښې لرې شوې؛ د
فرضونو سېټونه، د يا او شرط داخلول/ايستل او د
وروستي فرض ايستل د اصلي رسم په عين ترتيب دي۔
پايله او د لنډ نوم معنا عين ساتل شوي دي۔

\end{document}
